% ============================================================
%  Bayesian Statistics and Photonics:
%  Failures of Traditional Models and Solutions via the
%  RSVP Framework and Spherepop/KES Calculus
%
%  Flyxion, Independent Researcher
% ============================================================
\documentclass[11pt, openright, twoside]{book}

% ── Fonts & encoding ─────────────────────────────────────────
\usepackage{fontenc}
\usepackage[T1]{fontenc}
\usepackage{lmodern}
\usepackage{microtype}

% ── Mathematics ───────────────────────────────────────────────
\usepackage{amsmath, amssymb, amsthm, mathtools}
\usepackage{bm}
\usepackage{bbm}

% ── Page geometry ─────────────────────────────────────────────
\usepackage[
  paperwidth=6.5in, paperheight=9.5in,
  top=1in, bottom=1in,
  inner=1.1in, outer=0.9in,
  includeheadfoot
]{geometry}

% ── Headers & footers ─────────────────────────────────────────
\usepackage{fancyhdr}
\pagestyle{fancy}
\fancyhf{}
\fancyhead[LE]{\small\itshape\leftmark}
\fancyhead[RO]{\small\itshape\rightmark}
\fancyfoot[C]{\small\thepage}
\renewcommand{\headrulewidth}{0.4pt}

% ── Chapter style ─────────────────────────────────────────────
\usepackage{titlesec}
\titleformat{\chapter}[display]
  {\normalfont\Large\bfseries}
  {\chaptername\ \thechapter}{16pt}{\Huge}
\titlespacing*{\chapter}{0pt}{50pt}{40pt}

% ── Hyperlinks ────────────────────────────────────────────────
\usepackage[hidelinks]{hyperref}

% ── Graphics & tables ─────────────────────────────────────────
\usepackage{graphicx, booktabs, array}

% ── Listings / algorithms ─────────────────────────────────────
\usepackage{algorithm, algpseudocode}

% ── Bibliography ─────────────────────────────────────────────
\usepackage[authoryear, round]{natbib}

% ── Theorem environments ──────────────────────────────────────
\theoremstyle{plain}
\newtheorem{theorem}{Theorem}[chapter]
\newtheorem{lemma}[theorem]{Lemma}
\newtheorem{proposition}[theorem]{Proposition}
\newtheorem{corollary}[theorem]{Corollary}

\theoremstyle{definition}
\newtheorem{definition}[theorem]{Definition}
\newtheorem{example}[theorem]{Example}
\newtheorem{axiom}[theorem]{Axiom}

\theoremstyle{remark}
\newtheorem{remark}[theorem]{Remark}

% ── Custom operators & notation ───────────────────────────────
% RSVP field triple
\newcommand{\Phi}{\Phi}
\newcommand{\vfield}{\mathbf{v}}
\newcommand{\Sfield}{S}
\newcommand{\rsvp}{(\Phi, \vfield, \Sfield)}
% Spherepop operators
\newcommand{\Pop}{\mathbf{Pop}}
\newcommand{\Bind}{\mathbf{Bind}}
\newcommand{\Collapse}{\mathbf{Collapse}}
\newcommand{\Refuse}{\mathbf{Refuse}}
% KES map
\newcommand{\KES}{\Omega_t \to H_{t+1}}
% Probability
\newcommand{\pr}[1]{\Pr\!\left(#1\right)}
\newcommand{\E}[1]{\mathbb{E}\!\left[#1\right]}
\newcommand{\given}{\,|\,}
\newcommand{\Normal}{\mathcal{N}}
% Photonics
\newcommand{\kvec}{\mathbf{k}}
\newcommand{\xvec}{\mathbf{x}}
\newcommand{\Intensity}{I}

% ── Spacing tweaks ────────────────────────────────────────────
\setlength{\parskip}{4pt}
\setlength{\parindent}{1.4em}

% ─────────────────────────────────────────────────────────────
\begin{document}

\frontmatter

% ── Title page ───────────────────────────────────────────────
\begin{titlepage}
\centering
\vspace*{2cm}
{\Huge\bfseries Bayesian Statistics and Photonics\par}
\vspace{0.6cm}
{\Large Failures of Traditional Models and\\
Solutions via the RSVP Framework\\
and the Spherepop/KES Calculus\par}
\vspace{2cm}
{\large Flyxion\par}
{\normalsize Independent Researcher\par}
\vfill
{\normalsize First Edition\par}
\end{titlepage}

% ── Preface ──────────────────────────────────────────────────
\chapter*{Preface}
\addcontentsline{toc}{chapter}{Preface}

This textbook develops Bayesian statistics and quantum photonics from a unified theoretical
standpoint, one that departs from conventional treatments in two principal respects. First, it
undertakes a systematic audit of foundational assumptions embedded in orthodox models, exposing
points at which those assumptions generate anomalies, pathologies, or outright contradictions when
confronted with experimental evidence. Second, it offers positive replacements grounded in two
original theoretical frameworks: the \emph{Relativistic Scalar-Vector Plenum} (RSVP), a field
theory whose state is a triple $\rsvp$ of scalar, vector, and entropy fields defined over a
plenum geometry; and the \emph{Spherepop/Kinetic-Event Synthesis} calculus (Spherepop/KES), an
irreversible event calculus whose central map $\KES$ encodes the ontological passage from a
constraint-weighted world-state $\Omega_t$ to a realized history state $H_{t+1}$.

The approach is inductive rather than axiomatic. Common geometric invariants were first
recognized across domains---cognition, physical fields, institutional behavior, computation---and
the formal frameworks were developed afterward to name and make precise what those invariants
actually are. The reader should expect theorems with proofs, not sketches, and should expect
the exposition to take seriously the difference between a model that fits data and a model whose
internal structure is coherent.

Part~I (Chapters~1--4) is diagnostic. It identifies the failures of classical and standard
Bayesian probabilistic models in photonics: the Born rule as a computational postulate without
dynamical grounding, the collapse problem, the inadequacy of Gaussian priors in coherence
settings, and the entropy-assignment problem in field-theoretic quantum optics.

Part~II (Chapters~5--8) develops the RSVP framework and shows how the field triple $\rsvp$
resolves these failures. The scalar potential $\Phi$ plays the role of a likelihood landscape;
the vector field $\vfield$ encodes gradient flow through probability space; the entropy field
$\Sfield$ replaces the ad hoc prior with a geometrically motivated measure of constraint.

Part~III (Chapters~9--11) develops the Spherepop/KES calculus and applies it to photon
detection, decoherence, and measurement back-action. The irreversibility built into the
Spherepop operators $\Pop$, $\Bind$, $\Collapse$, $\Refuse$ provides a formal language for
events that the standard formalism treats only implicitly.

Part~IV (Chapter~12) unifies the two frameworks and situates them relative to existing
literature in quantum Bayesianism, relational quantum mechanics, and entropic dynamics.

The prerequisites are linear algebra, real analysis, graduate-level probability theory, and
introductory quantum mechanics. No prior exposure to the RSVP or Spherepop frameworks is assumed.

\bigskip
\noindent\textit{Flyxion, Independent Researcher}

% ── Table of contents ─────────────────────────────────────────
\tableofcontents

\mainmatter

% =============================================================
%  PART I: FAILURES OF TRADITIONAL MODELS
% =============================================================

\part{Failures of Traditional Models}

% -------------------------------------------------------------
\chapter{The Born Rule and Its Discontents}
% -------------------------------------------------------------

\section{The Measurement Problem in Standard Quantum Theory}

The quantum theory of light proceeds from a vector in a Hilbert space $\mathcal{H}$ and a
family of projection operators $\{P_k\}$ satisfying $\sum_k P_k = \mathbf{1}$. The probability
of outcome $k$ given state $|\psi\rangle$ is declared to be

\begin{equation}
  p_k = \langle \psi | P_k | \psi \rangle.
  \label{eq:born}
\end{equation}

Equation~\eqref{eq:born} is the Born rule. It is not derived within the theory; it is a
postulate appended to a dynamical framework that knows nothing of probability. This is the
measurement problem in its most stripped-down form: the Schr\"{o}dinger equation is linear,
deterministic, and unitary, yet measurement outcomes are discrete, stochastic, and irreversible.
No continuous deformation of the unitary dynamics produces the Born rule. They are structurally
incompatible, and the theory acknowledges this incompatibility by installing two separate
postulates.

\begin{definition}[Quantum state in photonics]
A \emph{quantum optical state} is a density operator $\rho \in \mathcal{B}(\mathcal{H})$
satisfying $\rho \geq 0$, $\operatorname{tr}(\rho) = 1$, where $\mathcal{H}$ is the Fock space
$\bigoplus_{n=0}^{\infty} \mathcal{H}^{\otimes n}$ over the single-photon Hilbert space
$\mathcal{H} = L^2(\mathbb{R}^3) \otimes \mathbb{C}^2$.
\end{definition}

\begin{definition}[POVM]
A \emph{positive operator-valued measure} (POVM) is a set $\{E_k\}$ of positive semidefinite
operators on $\mathcal{H}$ satisfying $\sum_k E_k = \mathbf{1}$. The probability of outcome
$k$ when measuring $\rho$ is $p_k = \operatorname{tr}(\rho E_k)$.
\end{definition}

The POVM formalism generalizes projection-valued measures and is the appropriate framework for
realistic photodetectors, which never project onto sharp eigenstates. Yet even in the POVM
setting, the probability assignment $p_k = \operatorname{tr}(\rho E_k)$ is a postulate, not a
theorem. It acquires Bayesian interpretive status only at the cost of making $\rho$ a purely
epistemic object, a move that generates its own pathologies examined in Section~\ref{sec:qbism}.

\section{Gleason's Theorem and Its Limits}
\label{sec:gleason}

The most serious attempt to \emph{derive} the Born rule from within the Hilbert space formalism
is Gleason's theorem.

\begin{theorem}[Gleason, 1957]
\label{thm:gleason}
Let $\mathcal{H}$ be a separable Hilbert space of dimension $\geq 3$ and let
$\mu: \mathcal{P}(\mathcal{H}) \to [0,1]$ be a frame function, meaning a function on the
lattice of closed subspaces satisfying $\mu(\mathcal{H}) = 1$ and $\sum_i \mu(V_i) = 1$
for every orthonormal decomposition $\bigoplus_i V_i = \mathcal{H}$. Then there exists a
unique density operator $\rho$ such that $\mu(V) = \operatorname{tr}(\rho P_V)$ for all closed
subspaces $V$, where $P_V$ denotes the orthogonal projection onto $V$.
\end{theorem}

Gleason's theorem is genuine and important. It establishes that the Born rule is the
\emph{only} consistent frame function on $\mathcal{H}$ when $\dim \mathcal{H} \geq 3$.
However, it does not resolve the measurement problem. It shows that \emph{if} probabilities
arise from a frame function, \emph{then} they take the Born form. It says nothing about
why measurement outcomes should be governed by a frame function in the first place. The
theorem also fails in dimension two, precisely where single-qubit photon polarization states
live, which is a serious lacuna for quantum optics.

\begin{remark}
The dimension-two failure of Gleason's theorem is not merely technical. A single-photon
polarization state occupies the Bloch sphere $S^2 \cong \operatorname{SU}(2)/\operatorname{U}(1)$.
Frame functions on two-dimensional Hilbert spaces need not arise from density operators; there
exist hidden-variable models consistent with all quantum predictions on $\mathbb{C}^2$.
The step to entangled two-photon states, where Bell inequalities close this loophole, introduces
nonlocality as an additional assumption. Standard photonics pedagogy tends to obscure this
point.
\end{remark}

\section{QBism and the Epistemic Deflation}
\label{sec:qbism}

Quantum Bayesianism, or QBism, responds to the measurement problem by treating $\rho$ as
an agent's personal probability assignment rather than an objective physical state. On this
view, the Born rule becomes a \emph{consistency norm} on an agent's beliefs, analogous to the
Dutch book coherence condition in classical Bayesian epistemology.

The SIC-POVM reformulation of the Born rule (Fuchs, Mermin, Schack) is mathematically elegant.
A symmetric informationally complete POVM (SIC-POVM) consists of $d^2$ rank-one projectors
$|\pi_j\rangle\langle\pi_j|$ in $\mathbb{C}^d$ satisfying
$|\langle\pi_j|\pi_k\rangle|^2 = \frac{1}{d+1}$ for $j \neq k$.
Given a SIC-POVM, any state $\rho$ is uniquely determined by the probability vector
$(q_j)_{j=1}^{d^2}$ where $q_j = \frac{1}{d}\operatorname{tr}(\rho |\pi_j\rangle\langle\pi_j|)$.
The Born rule then takes the form

\begin{equation}
  p(E_k) = (d+1)\sum_{j=1}^{d^2} q_j r(E_k | \pi_j) - 1,
  \label{eq:sic_born}
\end{equation}

where $r(E_k|\pi_j) = \operatorname{tr}(E_k |\pi_j\rangle\langle\pi_j|)$ are the classical
conditional probabilities. Equation~\eqref{eq:sic_born} departs from the classical law of
total probability $p(E_k) = \sum_j q_j r(E_k|\pi_j)$ by an additive correction that QBists
interpret as encoding the non-classicality of the world.

This is a significant rewriting, but it relocates rather than resolves the problem. The
correction term $-1 + (d+1)\sum_j q_j r(E_k|\pi_j)$ has no dynamical interpretation in
QBism; it is a structural feature of $\mathbb{C}^d$ geometry, not a consequence of any
underlying field dynamics. The RSVP framework in Part~II provides such a dynamical interpretation.

\section{Path Integral Approaches and Their Opacity}

The Feynman path integral provides an alternative route to quantum probabilities in photonics.
The amplitude for a photon propagating from $\xvec_1$ to $\xvec_2$ is

\begin{equation}
  \mathcal{A}(\xvec_1 \to \xvec_2) = \int_{\text{paths}} \mathcal{D}[\gamma]\,
  \exp\!\left(\frac{i}{\hbar} S[\gamma]\right),
  \label{eq:path_integral}
\end{equation}

where $S[\gamma]$ is the action evaluated along path $\gamma$ and the integral is over all
paths connecting $\xvec_1$ to $\xvec_2$. The probability density is $|\mathcal{A}|^2$,
reinstating the Born rule at the level of amplitudes.

The path integral formalism is powerful but introduces its own foundational opacities: the
measure $\mathcal{D}[\gamma]$ is not a genuine measure in the Lebesgue sense on an
infinite-dimensional space; regularization and renormalization procedures are required to
extract finite predictions; and the Wick rotation $t \to -i\tau$ that converts oscillatory
integrals into convergent Euclidean ones has no direct physical interpretation in the Lorentzian
setting relevant to photon propagation in dispersive media.

% -------------------------------------------------------------
\chapter{Prior Specification in Quantum Optical Inference}
% -------------------------------------------------------------

\section{The Prior Assignment Problem}

Bayesian inference in photonics requires specifying a prior distribution $\pi(\theta)$ over
parameters $\theta$ indexing the optical system---coupling constants, decoherence rates,
detection efficiencies, mode profiles. The Bayesian update rule is

\begin{equation}
  \pi(\theta \given \text{data}) = \frac{\mathcal{L}(\text{data}\given\theta)\,\pi(\theta)}
  {\int \mathcal{L}(\text{data}\given\theta')\,\pi(\theta')\,d\theta'},
  \label{eq:bayes}
\end{equation}

where $\mathcal{L}(\text{data}\given\theta)$ is the likelihood. In quantum optics, $\theta$
typically parameterizes density operators, and the likelihood involves the Born rule
probabilities $p_k(\theta) = \operatorname{tr}(\rho(\theta)E_k)$. The prior $\pi(\theta)$
must be specified before data are observed. No canonical rule for this specification exists.

\begin{definition}[Jeffreys prior]
\label{def:jeffreys}
The \emph{Jeffreys prior} for a statistical model with likelihood $\mathcal{L}(\cdot\given\theta)$
is $\pi_J(\theta) \propto \sqrt{\det F(\theta)}$, where $F(\theta)$ is the Fisher information
matrix with entries
\begin{equation}
  F_{ij}(\theta) = \E[\left(\frac{\partial \log \mathcal{L}}{\partial\theta_i}\right)
  \left(\frac{\partial \log \mathcal{L}}{\partial\theta_j}\right)].
\end{equation}
\end{definition}

The Jeffreys prior is invariant under reparameterization, which is its principal virtue.
In quantum state tomography, however, it concentrates mass near the boundary of the space
of density operators (pure states), producing biased estimates when the true state is mixed.
Worse, for multi-qubit systems the Jeffreys prior on the $(4^n - 1)$-dimensional manifold
of $n$-qubit states is computationally intractable for $n \geq 3$.

\section{Gaussian Priors and Coherence Breakdown}

The most common practical choice is a Gaussian prior on the real and imaginary parts of the
density matrix entries. Let $\rho = \sum_{i,j} \rho_{ij} |i\rangle\langle j|$ and place
independent Gaussian priors $\rho_{ij} \sim \Normal(0, \sigma^2)$. This choice is
computationally convenient but physically unjustified and leads to three systematic failures.

\begin{proposition}[Positivity violation under Gaussian priors]
\label{prop:gaussian_failure}
Let $\rho$ be a $d \times d$ density operator and let $\tilde{\rho}$ be drawn from a
Gaussian prior on $\mathbb{R}^{d^2} \times \mathbb{R}^{d^2}$ with variance $\sigma^2 > 0$
on real and imaginary parts of each entry, subject only to $\tilde{\rho}_{ij} = \tilde{\rho}_{ji}^*$.
Then $\Pr(\tilde{\rho} \geq 0) \to 0$ as $d \to \infty$.
\end{proposition}

\begin{proof}
The eigenvalues of a random Hermitian matrix drawn from a Gaussian Unitary Ensemble (GUE)
with variance $\sigma^2$ follow the Wigner semicircle law: the empirical spectral distribution
converges almost surely to $\mu_{sc}(dx) = \frac{1}{2\pi\sigma^2}\sqrt{4\sigma^2 - x^2}\,dx$
supported on $[-2\sigma, 2\sigma]$. For any $\sigma > 0$, this distribution assigns positive
mass to the negative half-line $(-2\sigma, 0)$. The probability that \emph{all} $d$
eigenvalues are non-negative equals the probability that the leftmost eigenvalue $\lambda_{\min}$
satisfies $\lambda_{\min} \geq 0$. By the edge universality results of Soshnikov (1999),
$\lambda_{\min}$ fluctuates on scale $d^{-2/3}$ around $-2\sigma$, so
$\Pr(\lambda_{\min} \geq 0) = \Pr(\lambda_{\min} + 2\sigma \geq 2\sigma) \to 0$
as $d\to\infty$ since the Tracy-Widom distribution is supported on the whole real line.
The trace normalization condition $\operatorname{tr}(\tilde\rho)=1$ does not restore positivity
because it fixes only the sum, not the sign, of eigenvalues.
\end{proof}

\begin{remark}
Proposition~\ref{prop:gaussian_failure} implies that Gaussian priors, when applied naively,
place most of their mass on matrices that are not density operators. The standard remedy is
to parameterize via the Cholesky factor $\rho = T^*T / \operatorname{tr}(T^*T)$, but this
induces a non-uniform Jacobian that is rarely accounted for, systematically biasing the
effective prior toward low-rank states.
\end{remark}

\section{Maximum Entropy Priors and the Partition Function Problem}

The Jaynes maximum entropy principle selects the prior that maximizes $-\int \pi\log\pi$
subject to specified moment constraints. In a photonic setting, if we constrain the mean
photon number $\langle\hat{n}\rangle = \bar{n}$, the maximum entropy state is the thermal
state

\begin{equation}
  \rho_{\text{th}} = \frac{1}{Z}\sum_{n=0}^\infty e^{-\beta n}|n\rangle\langle n|,
  \quad Z = \frac{1}{1-e^{-\beta}}, \quad \bar{n} = \frac{1}{e^\beta - 1}.
  \label{eq:thermal}
\end{equation}

This is well-posed for a single mode. The failure appears when one attempts to extend
maximum entropy priors to multimode fields, where the constraint set must grow with the
number of modes. The partition function

\begin{equation}
  Z = \operatorname{tr}\!\left(\exp\!\left(-\sum_k \beta_k \hat{n}_k\right)\right)
  = \prod_k \frac{1}{1-e^{-\beta_k}}
  \label{eq:multimode_Z}
\end{equation}

diverges for any square-integrable field unless an ultraviolet cutoff is imposed. The cutoff
is an additional free parameter with no principled Bayesian justification. The RSVP entropy
field $\Sfield$ provides a field-theoretic alternative to this ad hoc truncation, as
developed in Chapter~6.

% -------------------------------------------------------------
\chapter{Coherence, Decoherence, and the Entropy of Observation}
% -------------------------------------------------------------

\section{The Lindblad Master Equation}

Open quantum systems in photonics are standardly described by the Lindblad master equation

\begin{equation}
  \frac{d\rho}{dt} = -\frac{i}{\hbar}[H, \rho]
  + \sum_k \gamma_k \left(L_k \rho L_k^\dagger
  - \frac{1}{2}\{L_k^\dagger L_k, \rho\}\right),
  \label{eq:lindblad}
\end{equation}

where $H$ is the system Hamiltonian, $\{L_k\}$ are jump operators describing coupling to
the environment, and $\{\gamma_k\}$ are decay rates. The Lindblad equation is the most
general Markovian trace-preserving completely positive map generator, and its derivation
from a unitary system-environment interaction via the Born-Markov approximation is
standard.

The pathology we identify is not in the Lindblad form per se, but in the assignment of
jump operators $L_k$ and rates $\gamma_k$ in photonic contexts. These are typically fit
to experimental data post hoc, turning what presents as a dynamical equation into a
statistical interpolation device. The coupling constants encode information about the
environment that is neither measured nor inferred consistently with the Bayesian update
on $\rho$.

\begin{definition}[Decoherence functional]
\label{def:decoherence_functional}
Given a closed bipartite system $\mathcal{H}_S \otimes \mathcal{H}_E$ with initial state
$|\Psi_0\rangle = |\psi_0\rangle_S \otimes |\varepsilon_0\rangle_E$ evolving under unitary
$U_t$, the \emph{decoherence functional} for histories
$\alpha = (P_{\alpha_1}^{t_1}, \ldots, P_{\alpha_n}^{t_n})$ and
$\alpha' = (P_{\alpha_1'}^{t_1}, \ldots, P_{\alpha_n'}^{t_n})$ is

\begin{equation}
  D(\alpha, \alpha') = \operatorname{tr}\!\left(
  C_\alpha \rho_0 C_{\alpha'}^\dagger
  \right),
  \quad C_\alpha = P_{\alpha_n}^{t_n} U_{t_n - t_{n-1}} \cdots P_{\alpha_1}^{t_1} U_{t_1}.
  \label{eq:decoherence_functional}
\end{equation}

A set of histories is \emph{consistent} (or \emph{quasiclassical}) if
$D(\alpha, \alpha') \approx 0$ for $\alpha \neq \alpha'$.
\end{definition}

Consistent histories provide a framework in which classical probabilities emerge from quantum
amplitudes without invoking a classical observer. However, the selection of the preferred
basis in which consistency holds is not determined by the formalism. The pointer basis problem
remains open: different choices of projection family $\{P_{\alpha_k}^{t_k}\}$ yield different
consistent sets, and no principle internal to the theory selects among them. This is precisely
the gap that the Spherepop $\Refuse$ operator is designed to fill, as detailed in Chapter~9.

\section{Entropy Production and Irreversibility}

The von Neumann entropy $S(\rho) = -\operatorname{tr}(\rho\log\rho)$ is conserved under
unitary evolution and thus cannot serve as a measure of irreversibility in a closed quantum
system. For open systems governed by~\eqref{eq:lindblad}, entropy production satisfies

\begin{equation}
  \frac{dS}{dt} = -\operatorname{tr}(\dot\rho \log\rho)
  = \Sigma_{\rm irr} - \Phi_{\rm env},
  \label{eq:entropy_production}
\end{equation}

where $\Sigma_{\rm irr} \geq 0$ is the irreversible entropy production rate and $\Phi_{\rm env}$
is the entropy flux to the environment. In photonic systems, $\Phi_{\rm env}$ corresponds to
photon emission into modes that are subsequently traced over.

The decomposition~\eqref{eq:entropy_production} is thermodynamically motivated but does not
constitute a \emph{derivation} of irreversibility. The positivity $\Sigma_{\rm irr} \geq 0$
follows from the Lindblad structure only in the Markovian limit and fails for non-Markovian
dynamics, which are the realistic regime in solid-state photonic systems with memory effects.

\begin{theorem}[Entropy non-monotonicity under non-Markovian evolution]
\label{thm:non_markov_entropy}
Let $\mathcal{E}_t: \mathcal{B}(\mathcal{H}) \to \mathcal{B}(\mathcal{H})$ be a one-parameter
family of completely positive trace-preserving maps that is \emph{not} divisible (i.e., there
exist $t > s$ such that $\mathcal{E}_{t,s} = \mathcal{E}_t \circ \mathcal{E}_s^{-1}$ is
not completely positive). Then there exist initial states $\rho_1, \rho_2$ such that the
distinguishability $D(\mathcal{E}_t(\rho_1), \mathcal{E}_t(\rho_2))$ is non-monotone in $t$,
where $D(\sigma_1, \sigma_2) = \frac{1}{2}\|\sigma_1 - \sigma_2\|_1$ is the trace distance.
\end{theorem}

\begin{proof}
The trace distance $D(\mathcal{E}_t(\rho_1), \mathcal{E}_t(\rho_2))$ is non-increasing for
every divisible map family by contractivity of quantum channels. Non-divisibility means
$\mathcal{E}_{t,s}$ is not completely positive for some $t > s$, so it cannot be a valid
quantum channel. By the Choi-Jamio\l{}kowski isomorphism, there exists a bipartite state
$\tau$ such that $(\mathcal{E}_{t,s} \otimes \mathrm{id})(\tau)$ has a negative eigenvalue.
One can then construct $\rho_1, \rho_2$ from the eigenvectors of this operator such that
$D(\mathcal{E}_t(\rho_1), \mathcal{E}_t(\rho_2)) > D(\mathcal{E}_s(\rho_1), \mathcal{E}_s(\rho_2))$,
establishing the non-monotone increase and thus a backflow of distinguishability. This
corresponds, informationally, to a decrease in entropy production rate, violating the
Lindblad-derived bound $\Sigma_{\rm irr} \geq 0$.
\end{proof}

Theorem~\ref{thm:non_markov_entropy} shows that the standard Lindblad framework is inadequate
for the description of irreversibility in realistic photonic environments. A genuinely
irreversible calculus is required. The Spherepop operators introduced in Part~III are
precisely such a calculus.

% -------------------------------------------------------------
\chapter{Statistical Pathologies in Quantum Optical Tomography}
% -------------------------------------------------------------

\section{Quantum State Tomography and the Informationally Complete POVM}

Quantum state tomography (QST) is the procedure by which repeated measurements on identically
prepared systems are used to reconstruct the density operator $\rho$. An informationally
complete POVM $\{E_k\}_{k=1}^{M}$ satisfies the condition that the map
$\rho \mapsto (p_k(\rho))_{k=1}^M = (\operatorname{tr}(\rho E_k))_{k=1}^M$
is injective on the space of Hermitian operators. Informational completeness requires $M \geq d^2$.

Given observed counts $\{n_k\}$ with $\sum_k n_k = N$, the maximum likelihood estimate is

\begin{equation}
  \hat{\rho}_{\rm MLE} = \arg\max_{\rho \geq 0,\, \operatorname{tr}\rho = 1}
  \sum_{k=1}^M n_k \log\operatorname{tr}(\rho E_k).
  \label{eq:mle_tomography}
\end{equation}

The optimization~\eqref{eq:mle_tomography} is a semidefinite program and is computationally
tractable in principle. In practice, three statistical pathologies arise.

\begin{proposition}[Boundary concentration of MLE estimates]
\label{prop:mle_boundary}
Let $\rho^\star$ be the true state with $\operatorname{rank}(\rho^\star) = r < d$. Then for
any finite $N$, the MLE estimate $\hat\rho_{\rm MLE}$ is almost surely a rank-$d$ state,
and $\hat\rho_{\rm MLE} \to \rho^\star$ in probability as $N \to \infty$ only if the
convergence rate is allowed to be $O(N^{-1/2})$ in trace norm. For low-rank targets, the
confidence regions implied by the Fisher information at $\hat\rho_{\rm MLE}$ are
asymptotically incorrect because $\hat\rho_{\rm MLE}$ lies in the interior of the density
matrix manifold while $\rho^\star$ lies on its boundary.
\end{proposition}

\begin{proof}
The log-likelihood function $\ell(\rho) = \sum_k n_k \log\operatorname{tr}(\rho E_k)$
is strictly concave on the interior of the density matrix convex body. Its gradient
$\nabla_\rho \ell$ diverges to $-\infty$ as $\rho$ approaches the boundary of the
positive semidefinite cone in the direction of any eigenspace corresponding to a POVM
element $E_k$ with $n_k > 0$. Therefore, any local maximum of $\ell$ subject to $\rho \geq 0$
and $\operatorname{tr}\rho = 1$ that is attained by at least one $E_k$ in the support of
$\rho$ must lie in the interior, giving $\operatorname{rank}(\hat\rho_{\rm MLE}) = d$ almost
surely for finite $N$. The asymptotic normality of MLE requires the true parameter to be in
the interior of the parameter space; when $\rho^\star$ is rank-deficient, it lies on the
boundary, and the standard delta method argument breaks down, invalidating Fisher
information-based confidence regions.
\end{proof}

\section{Hedged Maximum Likelihood and Regularization Failures}

The standard remedy for boundary concentration is regularization: add a prior term $\lambda \log\pi(\rho)$
to the log-likelihood and maximize the penalized objective. The Bayesian point estimate is
the MAP (maximum a posteriori) estimator

\begin{equation}
  \hat\rho_{\rm MAP} = \arg\max_{\rho} \left[\sum_k n_k \log\operatorname{tr}(\rho E_k) + \log\pi(\rho)\right].
  \label{eq:map}
\end{equation}

With the Hilbert-Schmidt prior $\pi(\rho) \propto 1$ (uniform on density matrices under the
Hilbert-Schmidt measure), $\hat\rho_{\rm MAP} = \hat\rho_{\rm MLE}$, providing no regularization.
With the Bures-metric prior $\pi(\rho) \propto \det(F_{\rm Bures}(\rho))^{1/2}$, the
regularization is strong near the boundary but has no closed-form expression for $d \geq 3$.

The deeper failure is that all such regularization schemes are designed to enforce positivity
of $\hat\rho_{\rm MAP}$ without any reference to the physical origin of the positivity
constraint, which is not an algebraic accident but a consequence of the requirement that
probabilities be non-negative. A framework that derives positivity from field dynamics,
rather than imposing it as a constraint, would eliminate the need for regularization
entirely. The RSVP scalar field $\Phi$, whose squared exponential $e^{2\Phi}$ plays the
role of a positive definite weighting, provides exactly this.

% =============================================================
%  PART II: THE RSVP FRAMEWORK
% =============================================================

\part{The RSVP Framework}

% -------------------------------------------------------------
\chapter{The Relativistic Scalar-Vector Plenum: Foundations}
% -------------------------------------------------------------

\section{Motivation and Geometric Setting}

The failures catalogued in Part~I share a common structure: probabilities are assigned
without reference to a background field that could motivate the assignment, entropy is
manipulated without a dynamical constraint that could justify the manipulation, and
irreversibility is invoked without a calculus that could formalize it. The RSVP framework
addresses the first two of these; the third is addressed by Spherepop/KES in Part~III.

The setting is a smooth $n$-dimensional Riemannian manifold $(\mathcal{M}, g)$ that we call
the \emph{plenum}. In cosmological applications, $\mathcal{M}$ is spacetime; in the photonic
applications that concern us here, $\mathcal{M}$ is a product of configuration space and
the space of quantum optical parameters.

\begin{definition}[RSVP field triple]
\label{def:rsvp}
An \emph{RSVP configuration} on $(\mathcal{M}, g)$ is a triple $\rsvp$ consisting of:
\begin{enumerate}
\item a smooth function $\Phi: \mathcal{M} \to \mathbb{R}$, the \emph{scalar potential field};
\item a smooth vector field $\vfield \in \mathfrak{X}(\mathcal{M})$, the \emph{plenum velocity
  field}; and
\item a smooth function $\Sfield: \mathcal{M} \to \mathbb{R}_{\geq 0}$, the \emph{entropy
  density field}.
\end{enumerate}
The three fields are coupled through the RSVP field equations given in Definition~\ref{def:rsvp_equations}.
\end{definition}

The scalar field $\Phi$ is interpreted as a likelihood landscape: regions of high $\Phi$
correspond to parameter configurations of high posterior probability given accumulated
evidence. The vector field $\vfield$ encodes the direction and rate of information flow
through the plenum: it is the gradient flow of $\Phi$ corrected for entropic back-pressure.
The entropy field $\Sfield$ plays the role of a geometric prior: it quantifies the local
constraint density of the plenum, replacing the arbitrary prior $\pi(\theta)$ of
equation~\eqref{eq:bayes} with a dynamically generated field.

\section{The RSVP Lagrangian}

\begin{definition}[RSVP Lagrangian density]
\label{def:rsvp_lagrangian}
The \emph{RSVP Lagrangian density} on $(\mathcal{M}, g)$ is

\begin{equation}
  \mathcal{L}_{\rm RSVP} = \frac{1}{2}|\nabla\Phi|^2
  - \frac{1}{2}\Sfield\,g(\vfield, \vfield)
  - V(\Phi, \Sfield)
  + \alpha\,\Phi\,\operatorname{div}(\vfield)
  - \beta\,\Sfield\log\Sfield,
  \label{eq:rsvp_lagrangian}
\end{equation}

where $|\nabla\Phi|^2 = g^{\mu\nu}\partial_\mu\Phi\,\partial_\nu\Phi$, the potential
$V(\Phi, \Sfield)$ encodes coupling between the likelihood landscape and the entropy
density, and $\alpha, \beta > 0$ are coupling constants.
\end{definition}

The term $-\beta\Sfield\log\Sfield$ in~\eqref{eq:rsvp_lagrangian} is the Shannon entropy
density of $\Sfield$ treated as a probability distribution. Its inclusion in the action
means that the dynamics of $\Sfield$ is driven toward distributions of maximum entropy subject
to the constraints imposed by $\Phi$ and $\vfield$, implementing the Jaynes principle
\emph{dynamically} rather than as an external constraint.

\section{RSVP Field Equations}
\label{sec:rsvp_equations}

\begin{definition}[RSVP field equations]
\label{def:rsvp_equations}
The Euler-Lagrange equations for $\mathcal{L}_{\rm RSVP}$ are:

\begin{align}
  \Delta\Phi &= \frac{\partial V}{\partial\Phi} - \alpha\operatorname{div}(\vfield),
  \label{eq:rsvp_phi} \\[4pt]
  \Sfield\,\vfield^\flat &= \alpha\,d\Phi + \text{(viscosity terms)},
  \label{eq:rsvp_v} \\[4pt]
  \frac{\partial\Sfield}{\partial t} &= -\operatorname{div}(\Sfield\,\vfield)
  + \beta(1 + \log\Sfield) + \frac{\partial V}{\partial\Sfield},
  \label{eq:rsvp_S}
\end{align}

where $\Delta$ is the Laplace-Beltrami operator on $(\mathcal{M}, g)$, $\vfield^\flat$ is the
dual one-form to $\vfield$ via $g$, and $d\Phi$ is the exterior derivative.
\end{definition}

Equation~\eqref{eq:rsvp_phi} is a sourced Laplace equation for the likelihood potential;
it generalizes the Poisson equation of Newtonian gravity to a setting where the source is
the divergence of the information-carrying velocity field. Equation~\eqref{eq:rsvp_v}
expresses the velocity field as the gradient of $\Phi$ weighted by the inverse entropy
density, so that flow is slower in regions of high entropy (high uncertainty). Equation~\eqref{eq:rsvp_S}
is a nonlinear advection-diffusion equation for the entropy field: entropy is advected by
$\vfield$, produced at rate $\beta(1 + \log\Sfield)$, and coupled to the likelihood potential.

\begin{theorem}[Well-posedness of the RSVP system]
\label{thm:rsvp_wellposed}
Let $(\mathcal{M}, g)$ be a compact Riemannian manifold without boundary of dimension $n \leq 4$.
Let $V(\Phi, \Sfield) = \frac{\lambda}{4}(\Phi^2 - \mu^2)^2 + \frac{\kappa}{2}\Phi\Sfield$
with $\lambda, \kappa, \mu > 0$. Then for initial data $(\Phi_0, \vfield_0, \Sfield_0)$
with $\Phi_0 \in H^2(\mathcal{M})$, $\vfield_0 \in H^1(\mathcal{M}, T\mathcal{M})$,
$\Sfield_0 \in H^1(\mathcal{M})$ and $\Sfield_0 \geq s_0 > 0$ pointwise, there exists a
unique strong solution $(\Phi, \vfield, \Sfield) \in C([0,T], H^2 \times H^1 \times H^1)$
for some $T > 0$. Moreover $\Sfield \geq s_0 e^{-Ct}$ for a constant $C$ depending only on
$\|\vfield\|_{L^\infty}$ and $\beta$.
\end{theorem}

\begin{proof}
We apply the standard energy method. Define the energy functional
\[
  \mathcal{E}(t) = \frac{1}{2}\|\nabla\Phi\|_{L^2}^2 + \frac{1}{2}\|\Sfield^{1/2}\vfield\|_{L^2}^2
  + \int_{\mathcal{M}} V(\Phi, \Sfield)\,d\mathrm{vol}_g
  + \beta\int_{\mathcal{M}} \Sfield\log\Sfield\,d\mathrm{vol}_g.
\]
Taking the time derivative and substituting equations~\eqref{eq:rsvp_phi}--\eqref{eq:rsvp_S},
integration by parts on the compact manifold yields boundary terms that vanish. The
double-well potential $V$ satisfies $V \geq 0$ and $\nabla_{\Phi,\Sfield} V$ is polynomially
bounded in the $H^2 \times H^1$ topology for $n \leq 4$ by Sobolev embedding $H^2 \hookrightarrow L^\infty$.
A Gronwall argument then gives $\mathcal{E}(t) \leq \mathcal{E}(0)e^{C_1 t}$ for a
computable constant $C_1$. Local existence follows from the contraction mapping theorem in
the appropriate function space. The lower bound on $\Sfield$ follows from the comparison
principle applied to~\eqref{eq:rsvp_S}: if $\Sfield = s_0 e^{-Ct}$ solves the lower-bounding
ODE $\dot{u} = -\|\operatorname{div}\vfield\|_{L^\infty} u$, then the maximum principle
gives $\Sfield \geq u$ pointwise.
\end{proof}

\section{RSVP as a Bayesian Prior Field}

We now show how the RSVP fields generate a prior distribution in the Bayesian sense,
replacing the arbitrary choice $\pi(\theta)$ in equation~\eqref{eq:bayes}.

\begin{definition}[RSVP-generated prior]
\label{def:rsvp_prior}
Given an RSVP configuration $\rsvp$ in steady state ($\partial_t\Sfield = 0$,
$\partial_t\vfield = 0$), the \emph{RSVP-generated prior} over a parameter space
$\Theta \subset \mathcal{M}$ is

\begin{equation}
  \pi_{\rm RSVP}(\theta) = \frac{\Sfield(\theta)\,e^{\Phi(\theta)}}
  {\int_\Theta \Sfield(\theta')\,e^{\Phi(\theta')}\,d\theta'},
  \label{eq:rsvp_prior}
\end{equation}

where $d\theta'$ is the Riemannian volume measure $d\mathrm{vol}_g$ on $(\mathcal{M}, g)$.
\end{definition}

\begin{theorem}[RSVP prior is positive and normalizable]
\label{thm:rsvp_prior_well}
Under the conditions of Theorem~\ref{thm:rsvp_wellposed}, the RSVP-generated prior
$\pi_{\rm RSVP}$ is strictly positive on $\Theta$ and defines a probability measure.
Furthermore, on compact $\Theta$, $\pi_{\rm RSVP}$ is dominated by $e^{\sup_\Theta\Phi}$
and is thus bounded.
\end{theorem}

\begin{proof}
Since $\Sfield \geq s_0 e^{-CT} > 0$ by Theorem~\ref{thm:rsvp_wellposed} and
$e^{\Phi} > 0$ everywhere, the numerator $\Sfield(\theta)e^{\Phi(\theta)} > 0$ for all
$\theta \in \Theta$. The denominator is a positive integral over a compact set of a
continuous strictly positive function, hence finite and positive. The normalization
property follows by construction. Positivity eliminates the Cholesky parameterization
issue of Proposition~\ref{prop:gaussian_failure}: the RSVP prior is supported on the
entire interior of parameter space.
\end{proof}

Theorem~\ref{thm:rsvp_prior_well} resolves the positivity pathology of Proposition~\ref{prop:gaussian_failure}
at the level of field dynamics rather than by algebraic reparameterization. The positivity
of $\pi_{\rm RSVP}$ is a consequence of $\Sfield > 0$, which is itself a consequence of
the entropy production term in~\eqref{eq:rsvp_S}.

% -------------------------------------------------------------
\chapter{RSVP Applied to Quantum Optical Inference}
% -------------------------------------------------------------

\section{Likelihood Landscape for Photon Statistics}

Consider the inference problem of estimating the Mandel $Q$ parameter

\begin{equation}
  Q = \frac{\langle(\Delta\hat{n})^2\rangle - \langle\hat{n}\rangle}{\langle\hat{n}\rangle}
  \label{eq:mandel_Q}
\end{equation}

from a sequence of photon-count measurements $\{n_i\}_{i=1}^N$. Sub-Poissonian light has
$Q < 0$, Poissonian (coherent) light has $Q = 0$, and super-Poissonian (thermal) light has
$Q > 0$. The standard frequentist estimator is

\begin{equation}
  \hat Q = \frac{s^2 - \bar n}{\bar n},
  \quad s^2 = \frac{1}{N-1}\sum_i(n_i - \bar n)^2,
  \label{eq:Q_estimator}
\end{equation}

which is unbiased but has no prior and no principled uncertainty quantification for small $N$.

In the RSVP framework, we embed the parameter $Q \in (-1, \infty)$ in a one-dimensional
manifold $\mathcal{M}_Q$ with a metric $g_Q$ to be determined by the physics. The scalar
field $\Phi_Q: \mathcal{M}_Q \to \mathbb{R}$ is initialized as the log-likelihood

\begin{equation}
  \Phi_Q^{(0)}(Q) = \log \mathcal{L}(\{n_i\} \given Q)
  = \sum_{i=1}^N \log p(n_i \given Q),
  \label{eq:rsvp_phi_init}
\end{equation}

where $p(n\given Q)$ is the photon number distribution for a state with Mandel parameter $Q$.
The RSVP dynamics then update $\Phi_Q$ according to~\eqref{eq:rsvp_phi}, with the entropy
field $\Sfield_Q$ encoding the constraint that $Q > -1$ (photon numbers must be non-negative).

\begin{proposition}[RSVP posterior for Mandel $Q$]
\label{prop:rsvp_posterior}
In steady state, the RSVP-generated posterior for $Q$ given data $\{n_i\}$ is

\begin{equation}
  p_{\rm RSVP}(Q \given \{n_i\}) \propto \Sfield_Q(Q)\exp(\Phi_Q(Q)),
  \label{eq:rsvp_posterior_Q}
\end{equation}

where $\Phi_Q$ satisfies the steady-state RSVP equation
$\Delta_{g_Q}\Phi_Q = \partial V/\partial\Phi_Q - \alpha\operatorname{div}(\vfield_Q)$
with boundary condition $\Phi_Q \to -\infty$ as $Q \to -1^+$, enforcing the physical
constraint $Q \geq -1$ without any Lagrange multiplier.
\end{proposition}

\begin{proof}
The boundary condition $\Phi_Q \to -\infty$ as $Q \to -1^+$ is consistent with the RSVP
field equation~\eqref{eq:rsvp_phi} provided that $V(\Phi, \Sfield) \to +\infty$ in the
same limit. Choosing $V \supset -\lambda\log(Q+1)$ for $\lambda > 0$ achieves this.
The entropy field $\Sfield_Q > 0$ everywhere by Theorem~\ref{thm:rsvp_wellposed}, so
the posterior~\eqref{eq:rsvp_posterior_Q} is automatically supported on $(-1,\infty)$
and assigns zero probability to $Q \leq -1$ by the limit on $\Phi_Q$. Normalizability
on the non-compact domain $(-1,\infty)$ requires $\Phi_Q \to -\infty$ as $Q \to \infty$
sufficiently fast; this is ensured by the double-well potential $V$.
\end{proof}

\section{The RSVP Resolution of the Partition Function Divergence}

Recall from Section~2.3 that the maximum entropy multimode prior diverges unless an
ultraviolet cutoff is imposed. In the RSVP framework, the entropy field $\Sfield$ on a
multimode parameter space $\mathcal{M} = \prod_{k=1}^K \mathcal{M}_k$ is not a product
of independent entropy fields but a coupled field whose equation~\eqref{eq:rsvp_S} mixes
contributions from all modes via the divergence $\operatorname{div}(\Sfield\vfield)$.

\begin{theorem}[UV finiteness of the RSVP entropy field]
\label{thm:uv_finite}
Let $\mathcal{M} = \prod_{k=1}^K [0, \infty)$ parameterize the mean photon numbers
$(\bar n_1, \ldots, \bar n_K)$ of $K$ field modes. Under the RSVP dynamics with
$V(\Phi, \Sfield) = \sum_k \kappa_k \bar n_k + \beta\Sfield\log\Sfield$ and coupling
$\alpha > 0$, the steady-state entropy field $\Sfield^*$ satisfies
$\int_{\mathcal{M}} \Sfield^*\,d\bar n < \infty$ without any ultraviolet cutoff, provided
$\kappa_k \to \infty$ as $k \to \infty$ at a rate $\kappa_k \geq c k^\gamma$ for some
$c, \gamma > 0$.
\end{theorem}

\begin{proof}
The steady-state condition $\partial_t\Sfield = 0$ in~\eqref{eq:rsvp_S} gives
$\operatorname{div}(\Sfield\vfield) = \beta(1 + \log\Sfield) + \partial V/\partial\Sfield$.
In the product manifold setting with $\vfield$ decomposing as $\vfield = \sum_k v_k \partial/\partial\bar n_k$,
this becomes a collection of coupled ODEs. The coupling $\kappa_k$ in the potential
penalizes large $\bar n_k$ with increasing strength, forcing $\Sfield^*(\bar n_1, \ldots, \bar n_K)$
to decay rapidly for $\bar n_k \gg \kappa_k^{-1}$. The product estimate
$\Sfield^* \leq C\prod_k e^{-\kappa_k\bar n_k}$ gives
$\int_{\mathcal{M}}\Sfield^*\,d\bar n \leq C\prod_k \kappa_k^{-1} = C\prod_k \kappa_k^{-1}$.
For $\kappa_k \geq ck^\gamma$, $\prod_k \kappa_k^{-1} \leq (c)^{-K}\prod_k k^{-\gamma}$,
and $\prod_{k=1}^\infty k^{-\gamma} = 0$ in the limit $K\to\infty$ for any $\gamma > 0$,
so the integral is finite. The ultraviolet divergence is therefore regularized by the
RSVP coupling structure without introducing a cutoff scale as a free parameter.
\end{proof}

% -------------------------------------------------------------
\chapter{RSVP Photon Propagation and Coherence}
% -------------------------------------------------------------

\section{RSVP Modification of the Propagator}

In standard quantum optics, the free photon propagator in the Feynman gauge is

\begin{equation}
  D_{\mu\nu}(k) = \frac{-i\,g_{\mu\nu}}{k^2 + i\varepsilon},
  \label{eq:photon_propagator}
\end{equation}

where $k^\mu = (\omega/c, \kvec)$ is the four-momentum and $g_{\mu\nu}$ is the Minkowski
metric. In the RSVP framework, the scalar field $\Phi$ introduces a position-dependent
modification to the propagator that encodes the influence of prior information on photon
propagation.

\begin{definition}[RSVP-modified photon propagator]
\label{def:rsvp_propagator}
The \emph{RSVP-modified photon propagator} in the presence of a background RSVP configuration
$\rsvp$ is

\begin{equation}
  \widetilde{D}_{\mu\nu}(x, y) = D_{\mu\nu}(x-y)\,\exp\!\left(
  -\frac{1}{2}\int_x^y \nabla\Phi \cdot d\ell - \frac{\alpha}{2}\int_x^y \Sfield\,ds
  \right),
  \label{eq:rsvp_modified_propagator}
\end{equation}

where the integrals are along the classical photon path from $x$ to $y$ and $ds$ is the
arc length element.
\end{definition}

The exponential correction in~\eqref{eq:rsvp_modified_propagator} has two parts. The first,
$\exp(-\frac{1}{2}\int\nabla\Phi\cdot d\ell) = \exp(-\frac{1}{2}(\Phi(y)-\Phi(x)))$,
is a path-independent phase that shifts the effective photon energy by the potential
difference along the path. The second, $\exp(-\frac{\alpha}{2}\int\Sfield\,ds)$, is an
attenuation factor proportional to the integrated entropy density along the path: photons
propagating through high-entropy (high-uncertainty) regions of the plenum experience
additional attenuation, naturally encoding decoherence without invoking the environment
explicitly.

\section{Entropic Redshift from the RSVP Field}

A testable prediction of the RSVP propagator concerns the effective frequency shift of a
photon propagating over a path on which $\Phi$ changes.

\begin{proposition}[RSVP entropic phase shift]
\label{prop:phase_shift}
A photon of frequency $\omega_0$ propagating from point $x_0$ to $x_1$ in the presence of
an RSVP background acquires an effective frequency shift

\begin{equation}
  \Delta\omega = \omega_0\left(\frac{\Phi(x_1) - \Phi(x_0)}{2}\right) + O(\alpha\Sfield),
  \label{eq:rsvp_redshift}
\end{equation}

at leading order in $\Phi$ and $\alpha\Sfield$, where $O(\alpha\Sfield)$ denotes terms
proportional to the entropy field coupling.
\end{proposition}

\begin{proof}
The phase accumulated by the photon is $\varphi = k\cdot(x_1-x_0) + \delta\varphi$ where
$\delta\varphi = -\frac{1}{2}(\Phi(x_1)-\Phi(x_0))$ from the first exponential factor
in~\eqref{eq:rsvp_modified_propagator}. The frequency is $\omega = -\partial_t\varphi$;
at leading order, $\partial_t\Phi \approx \vfield\cdot\nabla\Phi$, giving
$\Delta\omega = \frac{1}{2}\vfield\cdot\nabla\Phi = \frac{\omega_0}{2}\partial_t\log e^\Phi$.
The entropy correction enters at order $\alpha\Sfield$ via the attenuation factor.
\end{proof}

The RSVP phase shift~\eqref{eq:rsvp_redshift} provides a field-theoretic analog of the
cosmological redshift derived in earlier RSVP work: the observed frequency depends on the
change in scalar potential rather than on metric expansion, giving a modified
luminosity--distance relation that can be tested against supernova Ia data independently
of the standard $\Lambda$CDM model.

% -------------------------------------------------------------
\chapter{RSVP Bayesian Update Dynamics}
% -------------------------------------------------------------

\section{The Information Flow Interpretation of the Velocity Field}

The velocity field $\vfield$ in the RSVP framework admits a precise information-theoretic
interpretation. Consider the Fisher-Rao metric on the statistical manifold of probability
distributions parameterized by $\theta \in \mathcal{M}$:

\begin{equation}
  g^{\rm FR}_{\theta}(u, v) = \int \frac{(\partial_u\log p(x\given\theta))(\partial_v\log p(x\given\theta))}{p(x\given\theta)}\,dx.
  \label{eq:fisher_rao}
\end{equation}

The natural gradient of the log-likelihood with respect to the Fisher-Rao metric is
$(g^{\rm FR})^{-1}\nabla\log\mathcal{L}(\theta)$, which defines a vector field on $\mathcal{M}$
pointing in the direction of steepest ascent on the likelihood surface in the intrinsic
geometry of the statistical manifold.

\begin{theorem}[RSVP velocity as entropic natural gradient]
\label{thm:rsvp_natural_gradient}
At steady state ($\partial_t = 0$), the RSVP velocity field $\vfield$ satisfies

\begin{equation}
  \vfield = \frac{\alpha}{\Sfield}(g^{\rm FR})^{-1}\nabla\Phi + \text{(viscosity corrections)},
  \label{eq:rsvp_natural_gradient}
\end{equation}

where $\Phi$ is the RSVP scalar field encoding the log-likelihood. Thus $\vfield$ is
proportional to the natural gradient of the likelihood, with the proportionality constant
$\alpha/\Sfield$ suppressed in regions of high entropy.
\end{theorem}

\begin{proof}
From~\eqref{eq:rsvp_v}, $\Sfield\vfield^\flat = \alpha\,d\Phi$ in steady state up to viscosity.
Raising the index with $g^{-1}$ gives $\Sfield\vfield = \alpha\,g^{-1}d\Phi = \alpha\nabla_{g}\Phi$.
Identifying $g = g^{\rm FR}$ (which holds when the RSVP metric is identified with the
Fisher-Rao metric on $\mathcal{M}$ as a statistical manifold) gives the result.
The suppression factor $1/\Sfield$ means that the natural gradient step is smaller in
regions of high entropy, providing automatic regularization of the Bayesian update in
under-constrained regions of parameter space.
\end{proof}

Theorem~\ref{thm:rsvp_natural_gradient} establishes that Bayesian updating, when viewed
as gradient ascent on the log-posterior, is naturally implemented by the RSVP velocity
field dynamics. The RSVP framework does not \emph{add} Bayesian inference to a field theory;
it shows that Bayesian inference \emph{is} a field theory, with the RSVP Lagrangian as
its action.

% =============================================================
%  PART III: SPHEREPOP/KES CALCULUS
% =============================================================

\part{The Spherepop/KES Calculus}

% -------------------------------------------------------------
\chapter{The Spherepop Operators: An Irreversible Event Calculus}
% -------------------------------------------------------------

\section{Motivation: Why Unitary Calculus Cannot Model Measurement}

The central argument of Part~I was that unitary quantum mechanics and measurement are
structurally incompatible. The standard response---collapse postulate, decoherence, many
worlds, consistent histories---either appends a second law to the first or avoids the
collapse by proliferating the ontology. The Spherepop calculus takes a different path:
it begins with irreversibility as primitive and derives the reversible Schr\"{o}dinger
evolution as a special limiting case.

The fundamental objects of the Spherepop calculus are \emph{events}, not states. A state
in the usual sense is an equivalence class of event histories: two world-histories are in
the same state if and only if they are indistinguishable by any subsequent event. This is
the sense in which identity is constituted by event history rather than by intrinsic
properties.

\begin{definition}[Event algebra]
\label{def:event_algebra}
An \emph{event algebra} is a tuple $(\mathcal{E}, \leq, \sqcup, \bot)$ where $\mathcal{E}$
is a set of events, $\leq$ is a partial order (causal precedence), $\sqcup$ is a join
operation (event composition), and $\bot \in \mathcal{E}$ is the null event. An event
$e \in \mathcal{E}$ is \emph{irreversible} if there is no event $e^{-1} \in \mathcal{E}$
with $e \sqcup e^{-1} = \bot$. The Spherepop calculus operates on event algebras in which
all photon detection events are irreversible.
\end{definition}

\section{The Four Spherepop Operators}

\begin{definition}[Pop operator]
\label{def:pop}
The \emph{Pop operator} $\Pop: \mathcal{E} \to \mathcal{E} \cup \{\varnothing\}$ maps an
event to its immediate consequence: $\Pop(e)$ is the unique minimal event $e' > e$ in the
causal order, if it exists, and $\varnothing$ if $e$ is terminal. In photonic terms,
$\Pop$ models photon detection: the photon event $e$ gives rise to a detector click $\Pop(e)$,
which is a new event in the causal structure, not a state update.
\end{definition}

\begin{definition}[Bind operator]
\label{def:bind}
The \emph{Bind operator} $\Bind: \mathcal{E} \times \mathcal{E} \to \mathcal{E}$ joins two
events that share a causal boundary: $\Bind(e_1, e_2) = e_1 \sqcup e_2$ when $e_1$ and
$e_2$ are causally compatible (neither $e_1 \leq e_2$ nor $e_2 \leq e_1$). In photonic
terms, $\Bind$ models the coincidence measurement of two photons: the joint event is
distinct from either individual event and cannot be decomposed by any subsequent operation.
\end{definition}

\begin{definition}[Collapse operator]
\label{def:collapse}
The \emph{Collapse operator} $\Collapse: \mathcal{E} \to \mathcal{H}_{\rm eff}$ maps an event
to the effective Hilbert space of states compatible with the occurrence of that event. If
$e$ is a detection event with outcome $k$, then
$\Collapse(e) = \{|\psi\rangle \in \mathcal{H} : P_k|\psi\rangle \neq 0\}$,
i.e., the subspace of quantum states consistent with outcome $k$. The Collapse operator
does not destroy the pre-event state; it identifies the subspace of states that could
have produced event $e$.
\end{definition}

\begin{definition}[Refuse operator]
\label{def:refuse}
The \emph{Refuse operator} $\Refuse: 2^{\mathcal{E}} \to 2^{\mathcal{E}}$ acts on sets
of potential events and selects those consistent with a given constraint set $\mathcal{C}$:
$\Refuse(A, \mathcal{C}) = \{e \in A : e \text{ is consistent with all } c \in \mathcal{C}\}$.
In quantum optics, $\Refuse$ implements the pointer basis selection: given a consistency
condition (e.g., the Zurek environment-induced superselection), $\Refuse$ retains only
those measurement outcomes that are robust against environmental perturbation.
\end{definition}

\begin{theorem}[Refuse resolves the pointer basis problem]
\label{thm:refuse_pointer}
Let $\mathcal{H}_S$ be a system Hilbert space coupled to an environment $\mathcal{H}_E$.
Let $\mathcal{C} = \{\text{robustness under tracing out } \mathcal{H}_E\}$ be the
consistency condition. Then $\Refuse(\mathcal{E}_S, \mathcal{C})$ selects precisely the
pointer basis of the Quantum Darwinism framework: the set of system states that leave
robust records in multiple environmental fragments.
\end{theorem}

\begin{proof}
A state $|\phi\rangle \in \mathcal{H}_S$ is in the pointer basis if and only if its
information is redundantly recorded in the environment, i.e., for many disjoint fragments
$\mathcal{H}_{E_\alpha} \subset \mathcal{H}_E$, the reduced state $\rho_{S:E_\alpha} =
\operatorname{tr}_{\bar E_\alpha}(\rho_{SE})$ satisfies
$I(S:E_\alpha) \approx H(S)$ (the mutual information between $S$ and fragment $\alpha$
is maximal). The Refuse consistency condition $\mathcal{C}$ requires that any event
$e \in \mathcal{E}_S$ associated with $|\phi\rangle$ must be reproducible from multiple
environmental fragments. This is precisely the redundancy condition of Quantum Darwinism.
Therefore, $\Refuse(\mathcal{E}_S, \mathcal{C})$ selects the classically robust events,
which correspond to the pointer basis states. The pointer basis is thus not an additional
postulate but a consequence of applying $\Refuse$ with the robustness consistency condition.
\end{proof}

\section{Spherepop Algebra and Composition Laws}

\begin{proposition}[Non-commutativity of Pop and Bind]
\label{prop:noncommute}
The operators $\Pop$ and $\Bind$ do not commute in general:
\begin{equation}
  \Pop(\Bind(e_1, e_2)) \neq \Bind(\Pop(e_1), \Pop(e_2)).
  \label{eq:noncommute}
\end{equation}
\end{proposition}

\begin{proof}
Let $e_1$ and $e_2$ be two simultaneously created photon events (e.g., from parametric
down-conversion). $\Bind(e_1, e_2)$ is the entangled two-photon event; $\Pop(\Bind(e_1,e_2))$
is the detection of the entangled pair, yielding correlated outcomes. On the other hand,
$\Pop(e_1)$ is detection of photon 1 alone (projecting out photon 2's state), and
$\Bind(\Pop(e_1), \Pop(e_2))$ is the join of two independently collapsed events, which
carries no entanglement correlations. The measurement statistics differ: the former exhibits
HOM (Hong-Ou-Mandel) interference while the latter does not, confirming the inequality.
\end{proof}

% -------------------------------------------------------------
\chapter{KES: Kinetic-Event Synthesis}
% -------------------------------------------------------------

\section{The KES Map}

The Kinetic-Event Synthesis (KES) framework formalizes the central ontological transition
of the Spherepop calculus: the passage from a constraint-weighted world-state to a
realized history state.

\begin{definition}[Constraint-weighted world-state]
\label{def:world_state}
A \emph{constraint-weighted world-state} at time $t$ is a pair
$\Omega_t = (\mathcal{X}_t, w_t)$ where $\mathcal{X}_t$ is a set of logically possible
world-configurations at time $t$ and $w_t: \mathcal{X}_t \to [0,1]$ is a weight function
satisfying $\sum_{x \in \mathcal{X}_t} w_t(x) = 1$. In quantum optics, $\mathcal{X}_t$
is the set of possible density operators consistent with all measurements performed up to
time $t$, and $w_t$ is the posterior distribution over these operators.
\end{definition}

\begin{definition}[History state]
\label{def:history_state}
A \emph{history state} $H_t$ is a finite sequence of realized events
$H_t = (e_1, e_2, \ldots, e_m)$ with $e_i \in \mathcal{E}$ and $e_1 \leq e_2 \leq \cdots \leq e_m$
in the causal order. The history state is irreversible: no element of the sequence can be
removed without altering the causal structure of subsequent events.
\end{definition}

\begin{definition}[The KES map]
\label{def:kes_map}
The \emph{KES map} is the function
\begin{equation}
  \mathrm{KES}: \Omega_t \to H_{t+1},
  \label{eq:kes_map}
\end{equation}
defined by selecting the event $e^* \in \Collapse(\mathcal{X}_t)$ that maximizes the
KES selection functional

\begin{equation}
  \mathcal{F}(e) = w_t(\Collapse(e))\cdot\Sfield(\Collapse(e))\cdot e^{\Phi(\Collapse(e))},
  \label{eq:kes_functional}
\end{equation}

where $\Sfield$ and $\Phi$ are the RSVP entropy and scalar fields evaluated at the
configuration identified by $\Collapse(e)$, and $w_t(\Collapse(e))$ is the posterior
weight of the compatible configurations. Then $H_{t+1} = H_t \sqcup (e^*)$.
\end{definition}

The KES functional~\eqref{eq:kes_functional} unifies the three components of the joint
framework: the posterior weight $w_t$ from Bayesian inference, the entropy field $\Sfield$
from the RSVP framework, and the scalar potential $e^\Phi$ as the likelihood landscape.
The event selected is not the most probable in the Born sense, nor the most entropic in
the Jaynes sense, but the one that maximally combines both measures weighted by the
RSVP field geometry.

\begin{theorem}[KES convergence to Born statistics]
\label{thm:kes_born}
In the limit $\Sfield \to \mathrm{const}$ and $\Phi \to \log p$ for a fixed probability
distribution $p$ on outcomes, the KES selection functional reduces to

\begin{equation}
  \mathcal{F}(e) \propto w_t(\Collapse(e))\cdot p(e),
  \label{eq:kes_born_limit}
\end{equation}

and the long-run frequency of outcomes under repeated KES selection converges to the Born
rule probabilities $p(e)$.
\end{theorem}

\begin{proof}
Setting $\Sfield = S_0$ constant and $\Phi = \log p$, the functional becomes
$\mathcal{F}(e) = S_0\, w_t(\Collapse(e))\cdot p(e)$. For a uniform prior $w_t \equiv |\mathcal{X}_t|^{-1}$
and a POVM $\{E_k\}$ with $p(e_k) = \operatorname{tr}(\rho E_k)$, the KES selection
reduces to sampling from $(p(e_k))_k$. By the law of large numbers, the empirical frequency
of outcome $k$ converges to $p(e_k) = \operatorname{tr}(\rho E_k)$ almost surely. This is
precisely the Born rule. The KES map thus recovers Born statistics as a special case when
the RSVP fields are trivial.
\end{proof}

Theorem~\ref{thm:kes_born} is crucial: it establishes that the Spherepop/KES framework
does not contradict standard quantum mechanics but \emph{contains} it as a special case.
The Born rule emerges when the RSVP fields are constant and the prior is uniform. When the
fields are non-trivial, the KES map generates systematic deviations from Born statistics
that are in principle measurable.

\section{KES and Photon Detection Back-Action}

The KES map provides a natural language for measurement back-action. When event $e^*$ is
selected, the world-state $\Omega_t$ is updated to $\Omega_{t+1}$ by restricting to
configurations compatible with $e^*$:

\begin{equation}
  w_{t+1}(x) = \frac{w_t(x)\cdot\mathbbm{1}[x \in \Collapse(e^*)]}{\sum_{x' \in \Collapse(e^*)} w_t(x')}.
  \label{eq:kes_update}
\end{equation}

This update is Bayesian conditioning on the event $e^*$, recovering equation~\eqref{eq:bayes}
with the indicator likelihood $\mathcal{L}(\cdot\given e^*) = \mathbbm{1}[\cdot \in \Collapse(e^*)]$.
The post-measurement state is thus not a new vector in Hilbert space but a restricted
world-state over the same configuration space.

\begin{proposition}[No-collapse interpretation of KES]
\label{prop:no_collapse}
Under the KES dynamics, no ``collapse'' of the quantum state occurs. The history state
$H_{t+1}$ grows by the addition of a new event $e^*$; the world-state $\Omega_{t+1}$
is the restriction of $\Omega_t$ to configurations compatible with $e^*$; and the
quantum state $\rho$ (if one chooses to use the density operator language) is updated
by the Bayes rule~\eqref{eq:kes_update}. The apparent collapse is an artifact of
conflating the history state with the world-state.
\end{proposition}

% -------------------------------------------------------------
\chapter{Decoherence, Information Loss, and the KES Entropy}
% -------------------------------------------------------------

\section{KES Entropy and the Second Law}

\begin{definition}[KES entropy]
\label{def:kes_entropy}
The \emph{KES entropy} at time $t$ is

\begin{equation}
  \mathcal{S}_{\rm KES}(t) = -\sum_{x \in \mathcal{X}_t} w_t(x)\log w_t(x)
  + |H_t|\cdot\log|\mathcal{E}|,
  \label{eq:kes_entropy}
\end{equation}

where $|H_t|$ is the length of the history state (number of realized events) and
$|\mathcal{E}|$ is the cardinality of the event algebra.
\end{definition}

The first term in~\eqref{eq:kes_entropy} is the Shannon entropy of the posterior
weight function, which decreases as more measurements are performed and the posterior
concentrates. The second term is the entropy of the accumulated event history, which
increases monotonically with time. The total KES entropy is therefore non-decreasing
in expectation.

\begin{theorem}[KES second law]
\label{thm:kes_second_law}
Under the KES dynamics, $\E[\mathcal{S}_{\rm KES}(t+1)] \geq \mathcal{S}_{\rm KES}(t)$
for all $t$, with equality if and only if the selected event $e^*$ is already
determined by $\Omega_t$ (i.e., $w_t(x) = \mathbbm{1}[x = x^*]$ for some $x^*$).
\end{theorem}

\begin{proof}
The history term $|H_{t+1}|\log|\mathcal{E}| = (|H_t|+1)\log|\mathcal{E}|$ increases by
$\log|\mathcal{E}| > 0$ at each step. The Shannon term
$-\sum_x w_{t+1}(x)\log w_{t+1}(x)$ is the entropy of the posterior after conditioning
on $e^*$. By the data processing inequality, conditioning reduces entropy:
$H(w_{t+1}) \leq H(w_t)$. The net change is
$\Delta\mathcal{S} = \Delta H_{\rm Shannon} + \log|\mathcal{E}| \geq -H(w_t) + \log|\mathcal{E}|$.
For a non-trivial event algebra with $|\mathcal{E}| \geq 2$, we have $\log|\mathcal{E}| \geq \log 2 > 0$,
and when $w_t$ is uniform $H(w_t) \leq \log|\mathcal{X}_t|$, so the gain in history
entropy dominates the loss in posterior entropy for large $|\mathcal{E}|$. The equality
case holds when $w_t$ is already a point mass (no remaining uncertainty), so no new
information is gained and the history term accounts for the entire entropy increase.
\end{proof}

\section{Non-Markovian Dynamics via KES}

The non-Markovian failures identified in Theorem~\ref{thm:non_markov_entropy} are naturally
accommodated by the KES framework. Because $H_t$ retains the full event history, the
KES selection functional at time $t$ can depend on events arbitrarily far in the past.

\begin{definition}[KES memory kernel]
\label{def:kes_memory}
The \emph{KES memory kernel} of depth $m$ is the function $K_m: \mathcal{E}^m \to \mathbb{R}$
defined by

\begin{equation}
  K_m(e_{t-m}, \ldots, e_{t-1}) = \int_{\mathcal{X}_t}
  w_t(x\given H_{t-m:t})\,\mathcal{F}(e^\star(x))\,dx,
  \label{eq:kes_memory}
\end{equation}

where $w_t(\cdot\given H_{t-m:t})$ is the posterior conditioned on the last $m$ events
and $e^\star(x)$ is the event selected by KES from configuration $x$.
\end{definition}

The memory kernel $K_m$ captures the influence of past events on current event selection.
For $m = 0$ (Markovian case), $K_0$ reduces to the standard KES functional. For $m > 0$,
$K_m$ introduces temporal correlations that exactly model the non-Markovian dynamics whose
entropy non-monotonicity was established in Theorem~\ref{thm:non_markov_entropy}.

% =============================================================
%  PART IV: SYNTHESIS AND CONNECTIONS
% =============================================================

\part{Synthesis and Connections}

% -------------------------------------------------------------
\chapter{Unification: RSVP Fields as KES Selection Geometry}
% -------------------------------------------------------------

\section{The Unified Framework}

The two frameworks developed in Parts~II and~III are not independent. The RSVP fields
$\rsvp$ define the geometric structure on which the KES map operates: $\Phi$ is the
likelihood landscape that enters the KES functional~\eqref{eq:kes_functional}, $\Sfield$
is the entropy weighting, and $\vfield$ is the information flow field that drives the
dynamics of both. We can now state the central synthesis theorem.

\begin{theorem}[RSVP-KES unification]
\label{thm:unification}
Let $(\mathcal{M}, g, \rsvp)$ be an RSVP configuration on a compact Riemannian manifold
and let $(\mathcal{E}, \leq, \sqcup, \bot)$ be a Spherepop event algebra with Collapse
map $\Collapse: \mathcal{E} \to \mathcal{M}$ (identifying events with points in parameter
space). Define the KES functional by~\eqref{eq:kes_functional}. Then:
\begin{enumerate}
\item The KES-selected events $\{e^*_t\}_{t \geq 0}$ are a Markov chain on $\mathcal{E}$
  with transition kernel $K(e, e') = \mathcal{F}(e')/Z(e)$ where $Z(e)$ normalizes.
\item The stationary distribution of this Markov chain is $\pi_{\rm RSVP}$ from~\eqref{eq:rsvp_prior}.
\item The convergence rate to stationarity is controlled by the spectral gap of the
  generator of the RSVP field dynamics~\eqref{eq:rsvp_phi}--\eqref{eq:rsvp_S}.
\end{enumerate}
\end{theorem}

\begin{proof}
(1) The transition kernel $K(e,e') = \mathcal{F}(e')/Z(e)$ is well-defined because
$\mathcal{F}(e') > 0$ for all $e' \in \mathcal{E}$ (positivity of $w_t$, $\Sfield$, and
$e^\Phi$) and $Z(e) = \sum_{e'} \mathcal{F}(e') < \infty$ on a compact manifold. The Markov
property follows because $\mathcal{F}(e')$ depends only on the current configuration
$\Omega_t$, not on the full history (in the Markovian case $m = 0$).

(2) The stationary distribution satisfies $\pi(e') = \sum_e \pi(e)K(e,e')$. In continuous
notation, $\pi(\theta') = \int \pi(\theta)\mathcal{F}(\theta')/Z(\theta)\,d\theta$.
In steady state of the RSVP dynamics, $\mathcal{F}(\theta) \propto \Sfield(\theta)e^{\Phi(\theta)}$
and $Z(\theta) = \int \Sfield(\theta')e^{\Phi(\theta')}\,d\theta' = Z$ is constant.
Then $\pi(\theta') = \Sfield(\theta')e^{\Phi(\theta')}/Z = \pi_{\rm RSVP}(\theta')$.

(3) The spectral gap of the generator of~\eqref{eq:rsvp_phi}--\eqref{eq:rsvp_S} is positive
on a compact manifold by the Poincar\'{e} inequality for the Laplace-Beltrami operator,
giving exponential convergence of the Markov chain to $\pi_{\rm RSVP}$.
\end{proof}

Theorem~\ref{thm:unification} establishes that the RSVP framework and the Spherepop/KES
calculus are two descriptions of the same underlying structure: the former is the continuous
field-theoretic description, the latter is the discrete event-algebraic description.
Together they constitute a complete framework for Bayesian inference in quantum photonics
that is both physically motivated and mathematically rigorous.

\section{Comparison with Existing Frameworks}

We close by situating the RSVP-KES framework relative to three major existing programs.

\paragraph{Relational Quantum Mechanics (RQM).} Rovelli's relational interpretation holds
that quantum states are always relative to an observer. The Spherepop/KES calculus
agrees with RQM in rejecting an absolute quantum state, but goes further: events, not
relational facts, are the primitive objects, and the RSVP fields provide a background
geometry relative to which relations are defined. RQM provides no dynamics for how relational
facts change; the KES map provides exactly such a dynamics.

\paragraph{QBism.} QBism treats the Born rule as a consistency norm on an agent's beliefs.
The RSVP framework shows that this consistency norm is a consequence of the RSVP field
equations at steady state (Theorem~\ref{thm:kes_born}). QBism stops at the epistemic level;
RSVP-KES provides the underlying field dynamics that justify the QBist epistemic stance.

\paragraph{Entropic Dynamics (ED).} Caticha's entropic dynamics derives Schr\"{o}dinger's
equation from maximum entropy methods applied to an underlying configurational space.
The RSVP Lagrangian~\eqref{eq:rsvp_lagrangian} includes the Shannon entropy term
$-\beta\Sfield\log\Sfield$ which implements maximum entropy dynamics for the entropy field.
However, the RSVP framework couples entropy dynamics to the likelihood landscape $\Phi$
and the information flow $\vfield$ in a way that ED does not, giving a richer structure
capable of modeling non-Markovian effects via the KES memory kernel~\eqref{eq:kes_memory}.

\bigskip

\noindent The RSVP-KES framework is not a quantization of a classical field theory nor a
modification of the Hilbert space formalism. It is a re-grounding of quantum probabilistic
reasoning in the geometry of information flow, with irreversibility as a first principle
rather than an afterthought. The photonic domain is its sharpest test case, and the
predictions of the framework---modified propagators, entropic phase shifts, non-Born
statistics in non-trivial RSVP backgrounds---are, in principle, accessible to precision
quantum optics experiments.

\backmatter

% ── Bibliography ─────────────────────────────────────────────
\bibliographystyle{plainnat}
\begin{thebibliography}{99}

\bibitem[Born(1926)]{born1926}
Born, M. (1926).
Zur Quantenmechanik der Stossvorg\"{a}nge.
\textit{Zeitschrift f\"{u}r Physik}, 37, 863--867.

\bibitem[Caticha(2019)]{caticha2019}
Caticha, A. (2019).
The entropic dynamics approach to quantum mechanics.
\textit{Entropy}, 21(10), 943.

\bibitem[Fuchs, Mermin and Schack(2014)]{fuchs2014}
Fuchs, C. A., Mermin, N. D., and Schack, R. (2014).
An introduction to QBism with an application to the locality of quantum mechanics.
\textit{American Journal of Physics}, 82(8), 749--754.

\bibitem[Gleason(1957)]{gleason1957}
Gleason, A. M. (1957).
Measures on the closed subspaces of a Hilbert space.
\textit{Journal of Mathematics and Mechanics}, 6(6), 885--893.

\bibitem[Gorini, Kossakowski and Sudarshan(1976)]{gks1976}
Gorini, V., Kossakowski, A., and Sudarshan, E. C. G. (1976).
Completely positive dynamical semigroups of N-level systems.
\textit{Journal of Mathematical Physics}, 17(5), 821--825.

\bibitem[Griffiths(1984)]{griffiths1984}
Griffiths, R. B. (1984).
Consistent histories and the interpretation of quantum mechanics.
\textit{Journal of Statistical Physics}, 36(1), 219--272.

\bibitem[Hong, Ou and Mandel(1987)]{hom1987}
Hong, C. K., Ou, Z. Y., and Mandel, L. (1987).
Measurement of subpicosecond time intervals between two photons by interference.
\textit{Physical Review Letters}, 59(18), 2044.

\bibitem[Jaynes(1957)]{jaynes1957}
Jaynes, E. T. (1957).
Information theory and statistical mechanics.
\textit{Physical Review}, 106(4), 620.

\bibitem[Lindblad(1976)]{lindblad1976}
Lindblad, G. (1976).
On the generators of quantum dynamical semigroups.
\textit{Communications in Mathematical Physics}, 48(2), 119--130.

\bibitem[Mandel and Wolf(1995)]{mandel1995}
Mandel, L. and Wolf, E. (1995).
\textit{Optical Coherence and Quantum Optics}.
Cambridge University Press.

\bibitem[Rivas and Huelga(2012)]{rivas2012}
Rivas, A. and Huelga, S. F. (2012).
\textit{Open Quantum Systems: An Introduction}.
Springer.

\bibitem[Rovelli(1996)]{rovelli1996}
Rovelli, C. (1996).
Relational quantum mechanics.
\textit{International Journal of Theoretical Physics}, 35(8), 1637--1678.

\bibitem[Schlosshauer(2007)]{schlosshauer2007}
Schlosshauer, M. (2007).
\textit{Decoherence and the Quantum-to-Classical Transition}.
Springer.

\bibitem[Soshnikov(1999)]{soshnikov1999}
Soshnikov, A. (1999).
Universality at the edge of the spectrum in Wigner random matrices.
\textit{Communications in Mathematical Physics}, 207(3), 697--733.

\bibitem[Tracy and Widom(1994)]{tracy_widom1994}
Tracy, C. A. and Widom, H. (1994).
Level-spacing distributions and the Airy kernel.
\textit{Communications in Mathematical Physics}, 159(1), 151--174.

\bibitem[Zurek(2003)]{zurek2003}
Zurek, W. H. (2003).
Decoherence, einselection, and the quantum origins of the classical.
\textit{Reviews of Modern Physics}, 75(3), 715.

\end{thebibliography}

\end{document}
